\section*{Overview}

Wavelet tree is a static range-query structure for arrays when the queries care about \emph{values inside a subarray}:
how many numbers are \(\le x\), what is the \(k\)-th smallest, how many times does a value occur, and similar
questions.

\section*{When to Use It}

Use it when:

\begin{itemize}
  \item the array is static or mostly static,
  \item queries are on subarrays,
  \item the answer depends on the multiset of values inside that subarray, not just on a sum or minimum.
\end{itemize}

\section*{Core Idea}

The tree recursively splits the value range \([lo, hi]\) by its midpoint. At each node, it stores a prefix-count array
telling how many elements of the current subsequence went to the left child.

That one prefix array is enough to map a subarray \([l, r]\) in the parent node into the corresponding subarray inside
either child.

\section*{Key Insight}

The structure does not partition by \emph{indices}; it partitions by \emph{value ranges} while preserving relative order
inside each node. That is why it can answer order-statistics questions without sorting every queried subarray.

\section*{Worked Problem}

\paragraph{Problem.}
Given a static array, answer queries of the form:

\begin{itemize}
  \item \(k\)-th smallest value in \([l, r]\),
  \item count of values \(\le x\) in \([l, r]\).
\end{itemize}

\paragraph{Why a wavelet tree fits.}
At each level, the query only needs to know how many of the range elements went left. That tells you whether the answer
stays in the left value half or the right value half.

\section*{Correctness Intuition}

Every node stores exactly the subsequence of values whose numeric range belongs to that node. The prefix counts map query
indices from parent order into child order, so recursion always follows the correct subset of positions.

\section*{Complexity Analysis}

For value range size \(\sigma\):

\begin{itemize}
  \item build: \(O(n \log \sigma)\),
  \item \(k\)-th smallest: \(O(\log \sigma)\),
  \item count \(\le x\): \(O(\log \sigma)\),
  \item memory: \(O(n \log \sigma)\).
\end{itemize}

\section*{Implementation}

The reference code supports:

\begin{itemize}
  \item \texttt{kth(l, r, k)},
  \item \texttt{lte(l, r, x)}.
\end{itemize}

The array is assumed to be 1-indexed for query arguments in the usual wavelet-tree style.

\section*{Common Pitfalls}

\begin{itemize}
  \item Mixing 0-indexed array storage with 1-indexed query formulas.
  \item Forgetting that the structure is static; arbitrary updates are not cheap.
  \item Building over a huge raw value domain instead of compressing when appropriate.
\end{itemize}

\section*{Variants / Extensions}

\begin{itemize}
  \item Frequency of one exact value.
  \item Range sum over values in a numeric interval with extra prefix data.
  \item Wavelet matrix, which stores the same idea in a flatter layout.
\end{itemize}

\section*{Practice Problems}

\begin{itemize}
  \item K-th smallest on subarrays.
  \item Count values below a threshold on subarrays.
  \item Static range median or quantile queries.
\end{itemize}

\section*{References}

\begin{itemize}
  \item \href{https://usaco.guide/adv/wavelet}{USACO Guide: Wavelet Tree}.
  \item \href{https://github.com/kth-competitive-programming/kactl}{KACTL}.
  \item \href{https://codeforces.com/catalog}{Codeforces Catalog}.
\end{itemize}
